%==============================================================================
%  Relazione di progetto - Miglioramento di TabPFN-2.5 tramite LoRA fine-tuning
%  Compilare con: pdflatex relazione_lora.tex  (due volte per i riferimenti)
%==============================================================================
\documentclass[11pt,a4paper]{article}

\usepackage[utf8]{inputenc}
\usepackage[T1]{fontenc}
\usepackage[italian]{babel}
\usepackage{lmodern}
\usepackage{microtype}
\usepackage[margin=2.5cm]{geometry}
\usepackage{amsmath,amssymb}
\usepackage{booktabs}
\usepackage{multirow}
\usepackage{graphicx}
\usepackage{xcolor}
\usepackage{listings}
\usepackage{caption}
\usepackage{enumitem}
\usepackage[hidelinks]{hyperref}

% --- Stile per i blocchi di codice ------------------------------------------
\definecolor{codegray}{gray}{0.95}
\definecolor{codekw}{rgb}{0.13,0.29,0.53}
\definecolor{codecmt}{rgb}{0.0,0.5,0.0}
\definecolor{codestr}{rgb}{0.58,0.0,0.0}
\lstset{
  basicstyle=\ttfamily\footnotesize,
  backgroundcolor=\color{codegray},
  keywordstyle=\color{codekw}\bfseries,
  commentstyle=\color{codecmt}\itshape,
  stringstyle=\color{codestr},
  numbers=left,
  numberstyle=\tiny\color{gray},
  frame=single,
  framesep=4pt,
  breaklines=true,
  showstringspaces=false,
  language=Python,
  captionpos=b,
}

\title{\textbf{Adattamento leggero di TabPFN-2.5 tramite LoRA fine-tuning}\\
\large Implementazione, esperimenti e analisi su classificazione binaria}
\author{Progetto di Machine Learning}
\date{\today}

\begin{document}
\maketitle

\begin{abstract}
\noindent
Questo lavoro studia se sia possibile migliorare le prestazioni del
\emph{tabular foundation model} TabPFN-2.5 su task di classificazione binaria
mediante un fine-tuning leggero basato su \emph{Low-Rank Adaptation} (LoRA),
senza riaddestrare i pesi principali del modello. Viene descritta
l'implementazione completa della pipeline (caricamento dati, metriche,
iniezione e addestramento degli adapter LoRA) e vengono condotti tre
esperimenti: (i)~la valutazione della generalizzazione di un adapter
addestrato su dataset medici e testato su dataset esterni; (ii)~un'ablation
sulla capacità degli adapter; (iii)~un'ablation sulla dimensione dei dati di
addestramento. I risultati mostrano che il LoRA \emph{non} migliora la
capacità discriminativa di TabPFN-2.5: aumentando di 4 volte la capacità degli
adapter o di circa 60 volte la quantità di dati, le metriche di discriminazione
(AUC-ROC, F1) restano invariate, e un addestramento prolungato tende a
degradarle. L'unico beneficio consistente, seppur modesto, riguarda la
calibrazione delle probabilità (Expected Calibration Error). Si tratta di un
risultato negativo ottenuto con metodologia rigorosa, che conferma quanto i
pesi pre-addestrati di TabPFN-2.5 siano vicini all'ottimo per questa classe di
problemi.
\end{abstract}

\tableofcontents
\bigskip

%==============================================================================
\section{Introduzione}
%==============================================================================

\subsection{Contesto: TabPFN-2.5}
TabPFN è un \emph{tabular foundation model} che affronta la classificazione su
dati tabulari mediante \emph{in-context learning}: invece di essere addestrato
su uno specifico dataset, il modello riceve in input l'intero training set come
``contesto'' e produce le predizioni per le righe di test in un singolo forward
pass, senza ottimizzazione dei pesi a tempo di utilizzo. I pesi del modello
sono pre-addestrati una volta per tutte su milioni di dataset sintetici. In
questo progetto si utilizza la versione \textbf{2.5} (non la v3, che è il
default del pacchetto), che supporta dataset fino a circa $50\,000$ campioni.

\subsection{Obiettivo}
L'obiettivo generale del progetto è migliorare le prestazioni di TabPFN-2.5 su
classificazione binaria combinando modifiche esterne al modello (preprocessing,
feature selection, ensembling, post-processing) con una modifica interna
leggera tramite LoRA fine-tuning, sotto il vincolo di \emph{non} riaddestrare i
pesi principali. La presente relazione documenta la parte relativa al
\textbf{LoRA fine-tuning}, che costituisce il primo blocco implementato e
sperimentato.

\subsection{Low-Rank Adaptation (LoRA)}
LoRA è una tecnica di \emph{parameter-efficient fine-tuning}: dato un layer
lineare con matrice di pesi $W \in \mathbb{R}^{d_{out}\times d_{in}}$, invece di
aggiornare $W$ si congela e si aggiunge un termine a basso rango
\begin{equation}
  y = Wx + \frac{\alpha}{r}\, B A\, x,
  \qquad A \in \mathbb{R}^{r\times d_{in}},\ B \in \mathbb{R}^{d_{out}\times r},
  \label{eq:lora}
\end{equation}
dove $r \ll \min(d_{in},d_{out})$ è il \emph{rango} e $\alpha$ un fattore di
scala. Si addestrano soltanto le matrici $A$ e $B$. All'inizializzazione si pone
$B = 0$, cosicché il termine aggiuntivo è nullo e il modello parte
numericamente identico a quello pre-addestrato; l'adattamento emerge solo
durante l'addestramento. Il numero di parametri allenabili passa così da
$d_{out}\,d_{in}$ a $r\,(d_{in}+d_{out})$, tipicamente una piccola frazione del
totale.

%==============================================================================
\section{Materiali e metodi}
%==============================================================================

\subsection{Dataset e preprocessing}
Sono stati utilizzati due gruppi disgiunti di dataset binari da OpenML
(Tabella~\ref{tab:dataset}): cinque dataset \emph{medici} per il fine-tuning e
quattro dataset per la \emph{valutazione} della generalizzazione, mai usati in
addestramento.

\begin{table}[h]
\centering
\caption{Dataset utilizzati (OpenML). I dataset di valutazione non sono mai
stati impiegati per l'addestramento degli adapter.}
\label{tab:dataset}
\small
\begin{tabular}{llrr}
\toprule
Ruolo & Dataset (OpenML ID) & Campioni & Feature \\
\midrule
Fine-tuning & \texttt{diabetes} (37)            & 768  & 8  \\
Fine-tuning & \texttt{breast\_cancer} (1510)    & 569  & 30 \\
Fine-tuning & \texttt{heart\_disease} (53)      & 270  & 13 \\
Fine-tuning & \texttt{chronic\_kidney} (42972)  & 400  & 24 \\
Fine-tuning & \texttt{hepatitis} (55)           & 155  & 19 \\
\midrule
Valutazione & \texttt{thyroid} (1000)           & 3772 & 28 \\
Valutazione & \texttt{adult} (1590)             & 48842& 14 \\
Valutazione & \texttt{credit\_g} (31)           & 1000 & 20 \\
Valutazione & \texttt{blood\_transfusion} (1464)& 748  & 4  \\
\bottomrule
\end{tabular}
\end{table}

\paragraph{Correzione degli identificativi.} Durante l'implementazione due ID
forniti nelle specifiche iniziali si sono rivelati errati e sono stati corretti
dopo verifica diretta su OpenML:
\begin{itemize}[nosep]
  \item \texttt{chronic\_kidney}: l'ID 40922 corrisponde in realtà al dataset
  \texttt{Run\_or\_walk\_information} (attività fisica, non pertinente); è stato
  sostituito con \textbf{42972} (\texttt{chronic-kidney-disease});
  \item \texttt{thyroid}: l'ID 40082 corrisponde a un dataset di chimica con 347
  classi (\texttt{QSAR-DATASET-\dots}); è stato sostituito con \textbf{1000}
  (\texttt{hypothyroid}), binario.
\end{itemize}

\paragraph{Pipeline di caricamento.} È stato implementato un modulo
(\texttt{utils/data\_loader.py}) che, per ciascun dataset: (1)~scarica i dati da
OpenML con caching locale; (2)~codifica le feature categoriche con
\texttt{OrdinalEncoder}, distinguendo numerico e categorico in modo robusto al
\emph{backend} dei tipi (in \texttt{pandas}~3.0 le stringhe hanno dtype
\texttt{str} Arrow, non \texttt{object}); (3)~imputa i valori mancanti con la
mediana (\texttt{SimpleImputer}); (4)~codifica il target in formato binario
$\{0,1\}$ con \texttt{LabelEncoder}, rimuovendo spazi/tabulazioni accidentali;
(5)~scarta colonne identificative (es. \texttt{id}) che introdurrebbero
\emph{data leakage}. La suddivisione in train/validation è stratificata.

\subsection{Metriche di valutazione}
Sono state implementate (modulo \texttt{evaluation/metrics.py}) quattro metriche
per la classificazione binaria, due di discriminazione e due di calibrazione:
\begin{itemize}[nosep]
  \item \textbf{AUC-ROC}: area sotto la curva ROC, capacità di discriminazione;
  \item \textbf{F1 macro}: media non pesata dell'F1 per classe, adatta a dataset
  sbilanciati;
  \item \textbf{Brier score}: $\frac{1}{N}\sum_i (p_i - y_i)^2$, errore quadratico
  medio delle probabilità (calibrazione);
  \item \textbf{Expected Calibration Error (ECE)}: suddividendo le probabilità in
  $M=10$ bin di uguale ampiezza,
  \begin{equation}
    \mathrm{ECE} = \sum_{b=1}^{M} \frac{|B_b|}{N}\,
    \bigl|\,\overline{p}(B_b) - \overline{y}(B_b)\,\bigr|,
  \end{equation}
  dove $\overline{p}(B_b)$ e $\overline{y}(B_b)$ sono la confidenza media e
  l'accuratezza media nel bin $b$.
\end{itemize}

\subsection{Architettura di TabPFN-2.5 e punti di iniezione}
\label{sec:arch}
L'architettura (file \texttt{architectures/tabpfn\_v2\_5.py}) è composta da una
pila di $L=24$ blocchi identici. Ogni blocco contiene due moduli di attenzione e
un MLP:
\begin{itemize}[nosep]
  \item \texttt{per\_sample\_attention\_between\_features}: attenzione tra le
  feature di una stessa riga;
  \item \texttt{per\_column\_attention\_between\_cells}: attenzione tra le righe
  di una stessa colonna (con \emph{multi-query attention} per le righe di test).
\end{itemize}
Entrambe ereditano da una classe \texttt{Attention} che definisce quattro
proiezioni lineari \emph{senza bias}: \texttt{q\_projection},
\texttt{k\_projection}, \texttt{v\_projection}, \texttt{out\_projection}, tutte
$\mathbb{R}^{192\times192}$ (dimensione di embedding $d=192$, $3$ teste,
\texttt{head\_dim}$=64$). Il modello completo conta circa $11{,}0$ milioni di
parametri.

I \textbf{punti di iniezione naturali} per LoRA sono dunque le proiezioni
dell'attenzione. Adottando per default il classico schema ``QV'' (solo
\texttt{q\_projection} e \texttt{v\_projection}) si ottengono
\[
  L \times 2~\text{(attenzioni)} \times 2~\text{(proiezioni)} = 24\times2\times2 = 96
  \quad\text{adapter}.
\]

\subsection{Implementazione del modulo LoRA}
Il modulo \texttt{models/lora.py} implementa l'Eq.~\eqref{eq:lora} come
\emph{wrapper} attorno a un \texttt{nn.Linear} esistente, indipendentemente da
TabPFN (e quindi testabile in isolamento). Il nucleo è riportato nel
Listing~\ref{lst:lora}.

\begin{lstlisting}[caption={Nucleo del wrapper \texttt{LoRALinear} (semplificato).},label=lst:lora]
class LoRALinear(nn.Module):
    def __init__(self, base_layer, r=8, alpha=16.0, dropout=0.0):
        super().__init__()
        self.base = base_layer                 # congelato
        for p in self.base.parameters():
            p.requires_grad = False
        self.scaling = alpha / r
        self.lora_A = nn.Parameter(torch.empty(r, base_layer.in_features))
        self.lora_B = nn.Parameter(torch.empty(base_layer.out_features, r))
        nn.init.kaiming_uniform_(self.lora_A, a=math.sqrt(5))
        nn.init.zeros_(self.lora_B)            # B=0 -> parte come il modello base

    def forward(self, x):
        lora = (self.lora_dropout(x) @ self.lora_A.t()) @ self.lora_B.t()
        return self.base(x) + self.scaling * lora
\end{lstlisting}

Sono inoltre fornite funzioni per: iniettare gli adapter individuando per nome i
\texttt{nn.Linear} bersaglio (\texttt{inject\_lora\_adapters}); congelare tutti i
parametri tranne quelli LoRA (\texttt{mark\_only\_lora\_as\_trainable});
salvare/caricare i soli pesi degli adapter (pochi MB); e fondere gli adapter nei
pesi base per l'inferenza (\texttt{merge}). Il modulo è stato validato con test
unitari su CPU che verificano, fra l'altro, l'invarianza dell'output
all'inizializzazione ($\lVert y_{\text{base}}-y_{\text{LoRA}}\rVert_\infty = 0$)
e la corretta riproducibilità dopo salvataggio e ricarica degli adapter.

\subsection{Procedura di addestramento}
L'addestramento (modulo \texttt{models/tabpfn\_lora.py}) adotta un \emph{loop
custom e trasparente}, che riutilizza le primitive di preprocessing di TabPFN ma
controlla esplicitamente quali parametri vengono ottimizzati. Per ogni
\emph{meta-batch} (corrispondente a un dataset, suddiviso in contesto e query):
\begin{enumerate}[nosep]
  \item si imposta il contesto in-context (\texttt{fit\_from\_preprocessed});
  \item si calcolano i \emph{logits} grezzi sulle query
  (\texttt{forward(\dots, return\_raw\_logits=True)});
  \item si calcola la cross-entropy tra logits e label delle query;
  \item si esegue \texttt{backward} e un passo di ottimizzazione,
  \emph{aggiornando soltanto i parametri LoRA}.
\end{enumerate}
L'ottimizzatore è AdamW con \emph{learning rate} $10^{-4}$ e \emph{gradient
clipping}. Per evitare l'esaurimento della memoria GPU su contesti ampi si
abilita l'\emph{activation checkpointing} (ricalcolo delle attivazioni in fase di
backward anziché memorizzazione).

\paragraph{Early stopping.} Una porzione del training set (separata dal
validation usato per il confronto finale, per evitare \emph{leakage}) viene
riservata alla validazione interna: a ogni epoca si misura l'AUC, si conservano i
pesi LoRA migliori e ci si arresta dopo un numero prefissato di epoche senza
miglioramento, ripristinando il miglior checkpoint. Questo accorgimento si è
rivelato essenziale (cfr.\ Sez.~\ref{sec:datasize}).

%==============================================================================
\section{Esperimenti e risultati}
%==============================================================================

Tutti gli esperimenti sono stati eseguiti su GPU (Google Colab, NVIDIA T4).
Baseline e modello LoRA condividono lo stesso numero di estimatori in inferenza
per garantire un confronto equo.

\subsection{Esperimento 1: generalizzazione (training multi-dataset)}
\label{sec:exp5}
Un \emph{unico} adapter LoRA ($r=8$, $\alpha=16$, schema QV; $294{,}912$
parametri allenabili pari al $2{,}68\%$ del totale) è stato addestrato
congiuntamente sui cinque dataset medici e poi valutato sui quattro dataset di
test mai visti, per misurare la generalizzazione dell'adattamento. Durante
l'addestramento la media dell'AUC di validazione è salita da $0{,}920$ a
$0{,}930$, indicando che gli adapter \emph{apprendono} qualcosa di specifico al
dominio medico. I risultati sui dataset di test sono riportati in
Tabella~\ref{tab:exp5}.

\begin{table}[h]
\centering
\caption{Esperimento 1 -- TabPFN base vs.\ TabPFN+LoRA sui dataset di
valutazione. In \textbf{grassetto} il valore migliore per metrica e dataset.
Frecce: $\uparrow$ meglio se alto, $\downarrow$ meglio se basso.}
\label{tab:exp5}
\small
\begin{tabular}{llcccc}
\toprule
Dataset & Modello & AUC-ROC$\uparrow$ & F1$\uparrow$ & Brier$\downarrow$ & ECE$\downarrow$ \\
\midrule
\multirow{2}{*}{\texttt{adult}}
 & base & \textbf{0.9220} & 0.8056 & \textbf{0.0920} & 0.0116 \\
 & LoRA & 0.9219 & \textbf{0.8060} & 0.0921 & \textbf{0.0107} \\
\addlinespace
\multirow{2}{*}{\texttt{blood\_transfusion}}
 & base & \textbf{0.7908} & \textbf{0.6643} & \textbf{0.1419} & 0.0948 \\
 & LoRA & 0.7906 & \textbf{0.6643} & \textbf{0.1419} & \textbf{0.0909} \\
\addlinespace
\multirow{2}{*}{\texttt{credit\_g}}
 & base & 0.7750 & \textbf{0.6796} & 0.1677 & \textbf{0.0399} \\
 & LoRA & \textbf{0.7759} & \textbf{0.6796} & \textbf{0.1676} & 0.0454 \\
\addlinespace
\multirow{2}{*}{\texttt{thyroid}}
 & base & \textbf{1.0000} & \textbf{0.9953} & \textbf{0.0014} & \textbf{0.0038} \\
 & LoRA & \textbf{1.0000} & \textbf{0.9953} & 0.0016 & 0.0041 \\
\bottomrule
\end{tabular}
\end{table}

Le differenze sono tutte dell'ordine della terza--quarta cifra decimale, ovvero
entro il rumore statistico. Il LoRA \emph{non} migliora né peggiora in modo
significativo la discriminazione; si osserva un debole e non robusto vantaggio
sulla calibrazione (ECE migliore in 2 dataset su 4). Da notare che il dataset
medico tra i test (\texttt{thyroid}) è saturo (AUC $=1{,}0000$), quindi non offre
margine di osservazione del transfer in-dominio.

\subsection{Esperimento 2: ablation sulla capacità}
\label{sec:capacity}
Per verificare se l'assenza di miglioramento dipenda da una capacità
insufficiente degli adapter, l'Esperimento~1 è stato ripetuto con tre
configurazioni crescenti (Tabella~\ref{tab:capacity}).

\begin{table}[h]
\centering
\caption{Esperimento 2 -- configurazioni LoRA a capacità crescente. La migliore
AUC di validazione (sui dati medici) resta costante.}
\label{tab:capacity}
\small
\begin{tabular}{lrrc}
\toprule
Config. & Target & Param.\ allenabili & Best val AUC \\
\midrule
\texttt{r8\_qv}    & q, v       & $294{,}912$\ ($2{,}7\%$)  & 0.9294 \\
\texttt{r16\_qv}   & q, v       & $589{,}824$\ ($5{,}2\%$)  & 0.9295 \\
\texttt{r16\_qkvo} & q, k, v, o & $1{,}179{,}648$\ ($9{,}9\%$) & 0.9295 \\
\bottomrule
\end{tabular}
\end{table}

Quadruplicando i parametri allenabili (dal $2{,}7\%$ al $9{,}9\%$) la migliore
AUC di validazione resta sostanzialmente invariata ($0{,}9295$) e le metriche sui
dataset di test sono indistinguibili; in un caso (F1 su \texttt{credit\_g}) la
configurazione più capiente peggiora leggermente, suggerendo un principio di
overfitting. Si conclude che \textbf{la capacità degli adapter non è il fattore
limitante}.

\subsection{Esperimento 3: ablation sulla dimensione dei dati}
\label{sec:datasize}
La seconda ipotesi è che i dataset di fine-tuning, essendo piccoli
($155$--$768$ righe), forniscano troppi pochi passi di gradiente. Per testarla,
un adapter ($r=8$, QV) è stato addestrato e valutato \emph{in-distribution} su un
dataset medico ampio, \texttt{Cardiovascular-Disease} (OpenML~45547), ridotto a
$40\,000$ campioni per rispettare il limite di TabPFN ($\sim$$50\,000$):
$32\,000$ righe di training (circa $60$ volte i dataset piccoli, suddivise in
$\sim$$8$ chunk per epoca) e $8\,000$ di test.

\begin{table}[h]
\centering
\caption{Esperimento 3 -- TabPFN base vs.\ LoRA su \texttt{cardiovascular}
($32\,000$ righe di training, valutazione in-distribution).}
\label{tab:datasize}
\small
\begin{tabular}{lcccc}
\toprule
Modello & AUC-ROC$\uparrow$ & F1$\uparrow$ & Brier$\downarrow$ & ECE$\downarrow$ \\
\midrule
base & \textbf{0.7984} & \textbf{0.7291} & 0.1826 & 0.0275 \\
LoRA & 0.7983 & 0.7290 & \textbf{0.1824} & \textbf{0.0229} \\
\bottomrule
\end{tabular}
\end{table}

Nonostante $60$ volte più dati, il LoRA resta equivalente al baseline sulla
discriminazione (Tabella~\ref{tab:datasize}); l'unico miglioramento netto è di
nuovo sulla calibrazione (ECE $0{,}0229$ vs.\ $0{,}0275$). \textbf{Anche
l'ipotesi della scarsità di dati è quindi rigettata.}

Particolarmente istruttiva è la dinamica dell'addestramento
(Figura~\ref{fig:valauc}): la AUC di validazione raggiunge il massimo
($0{,}8045$) già all'epoca~5 e poi \emph{decresce} monotonicamente
($0{,}7789$ all'epoca~30), mentre la loss di training cresce. In altre parole, il
modello parte $\approx$ base (per via di $B=0$) e ogni ulteriore addestramento lo
\emph{allontana} dalla soluzione ottimale, peggiorandolo; solo l'early stopping
preserva le prestazioni. Si noti infine che \texttt{cardiovascular} non è saturo
(AUC $\approx 0{,}80$): esisteva margine di miglioramento, eppure il LoRA non lo
ha colto.

\begin{figure}[h]
\centering
\begin{tabular}{c}
\fbox{\parbox{0.85\textwidth}{\centering\footnotesize
\textit{[Andamento della AUC di validazione per epoca su
\texttt{cardiovascular}.]}\\[2pt]
Ep.\ 5: \textbf{0.8045} (max) \;$\to$\; Ep.\ 15: 0.8011 \;$\to$\;
Ep.\ 20: 0.7966 \;$\to$\; Ep.\ 25: 0.7882 \;$\to$\; Ep.\ 30: 0.7789\\[2pt]
La AUC decresce dopo l'epoca 5: l'addestramento prolungato degrada il modello.}}
\end{tabular}
\caption{Dinamica dell'early stopping nell'Esperimento~3. Il valore migliore è
all'epoca 5; oltre, le prestazioni peggiorano.}
\label{fig:valauc}
\end{figure}

%==============================================================================
\section{Discussione}
%==============================================================================
I tre esperimenti convergono verso una conclusione coerente
(Tabella~\ref{tab:summary}). Il LoRA fine-tuning non migliora la capacità
discriminativa di TabPFN-2.5 su classificazione binaria tabulare. Le due
spiegazioni alternative più naturali sono state escluse sperimentalmente:
\begin{itemize}[nosep]
  \item \textbf{non è un problema di capacità}: un aumento di $4\times$ dei
  parametri allenabili non produce alcun cambiamento (Sez.~\ref{sec:capacity});
  \item \textbf{non è un problema di dati}: un aumento di $\sim$$60\times$ dei
  dati di training non produce miglioramenti, e un addestramento prolungato
  degrada il modello (Sez.~\ref{sec:datasize}).
\end{itemize}

\begin{table}[h]
\centering
\caption{Sintesi degli esperimenti e delle ipotesi testate.}
\label{tab:summary}
\small
\begin{tabular}{lll}
\toprule
Ipotesi & Esperimento & Esito \\
\midrule
Generalizzazione fuori dominio & Exp.~1 (medici $\to$ test) & LoRA $\approx$ base, non distruttivo \\
Capacità insufficiente & Exp.~2 (ablation $r$, target) & nessun cambiamento \\
Dati insufficienti & Exp.~3 ($60\times$ dati) & nessun gain; over-training degrada \\
\bottomrule
\end{tabular}
\end{table}

La spiegazione residua, e più solida, è che \textbf{i pesi pre-addestrati di
TabPFN-2.5 sono già vicini all'ottimo} per questa classe di problemi: il suo
in-context learning estrae quasi tutto il segnale disponibile, e un adattamento a
basso rango non può aggiungervi capacità discriminativa, potendo solo
allontanarsi dal buon punto di partenza. Un aspetto positivo, da non trascurare,
è che l'adattamento risulta \emph{non distruttivo}: pur essendo addestrato solo
su dati medici, non degrada le prestazioni sui dataset di test non medici,
evidenziando robustezza allo \emph{shift} di dominio. Inoltre, l'unico segnale
ricorrente -- per quanto modesto -- è il miglioramento della \emph{calibrazione}
(ECE), il che suggerisce che eventuali benefici vadano cercati nel
post-processing piuttosto che nell'adattamento dei pesi.

\paragraph{Note implementative rilevanti.} Durante lo sviluppo sono emersi alcuni
accorgimenti tecnici necessari: (i)~l'early stopping è indispensabile, poiché
l'addestramento può degradare il modello; (ii)~su dataset grandi occorre
abilitare l'\emph{activation checkpointing} per non esaurire la memoria della GPU
in fase di training; (iii)~per dataset oltre il limite ufficiale va impostato
\texttt{ignore\_pretraining\_limits=True}.

%==============================================================================
\section{Conclusioni e lavori futuri}
%==============================================================================
È stato implementato e validato un sistema completo per il LoRA fine-tuning di
TabPFN-2.5: caricamento e preprocessing dei dati, quattro metriche di
valutazione, iniezione e addestramento degli adapter, ed early stopping. Sul
piano sperimentale, mediante tre esperimenti con relative ablation, si è stabilito
con metodologia rigorosa che il LoRA non migliora la discriminazione di
TabPFN-2.5 su classificazione binaria, e che ciò non dipende né dalla capacità
degli adapter né dalla quantità di dati, bensì dalla prossimità all'ottimo dei
pesi pre-addestrati. Il risultato, sebbene negativo, è informativo e
metodologicamente solido.

Il debole ma consistente miglioramento sulla calibrazione indica la direzione
naturale per il proseguimento del progetto: la \textbf{pipeline esterna}
(preprocessing adattivo, feature selection, ensembling e, soprattutto,
\emph{post-processing} con calibrazione e ottimizzazione della soglia per l'F1),
dove ci si attende un segnale più marcato, in particolare sui dataset
sbilanciati.

%==============================================================================
\appendix
\section{Struttura del codice}
%==============================================================================
\begin{itemize}[nosep]
  \item \texttt{utils/data\_loader.py} -- caricamento, caching, encoding e split
  dei dataset OpenML.
  \item \texttt{evaluation/metrics.py} -- AUC-ROC, F1 macro, Brier, ECE; tabelle
  comparative e salvataggio CSV.
  \item \texttt{models/lora.py} -- implementazione di LoRA (\texttt{LoRALinear},
  iniezione, salvataggio/caricamento, merge).
  \item \texttt{models/tabpfn\_lora.py} -- integrazione con TabPFN: costruzione
  del classificatore con adapter, loop di training, early stopping, esperimenti
  (\texttt{run\_lora\_exp5}, \texttt{run\_lora\_ablation},
  \texttt{run\_lora\_datasize\_experiment}).
\end{itemize}

\noindent\textit{Codice e configurazione completi sono versionati nel repository
del progetto. Il modello TabPFN-2.5 è eseguito localmente; il primo download dei
pesi richiede un token di licenza Prior Labs.}

\end{document}
